\documentclass[a4paper,12pt]{dinbrief} 
\usepackage{ngerman}
\usepackage[utf8]{inputenc}
%\usepackage{eurosym}
%\usepackage[top=4cm,left=2.9cm,right=3cm,bottom=4cm]{geometry}


\begin{document} 
\setbottomtexttop{2cm}
%\nowindowtics
%\nowindowrules
\subject{\textbf{Widerspruch Bescheid 4702-2401 vom 12.05.2026} }
\signature{Dr.~Fabian Schury}
\place{Aurachtal}
\Datum{\today} 
\backaddress{Dr.~Fabian Schury $|$ Reichenfelser Straße 10 $|$ 91086 Aurachtal}
\address{Dr.~Fabian Schury \\ Reichenfelser Straße 10 \\ 91086 Aurachtal}
 
\begin{letter} 
{   Zentrum Bayern Familie und Soziales\\
    Scanstelle Regensburg \\[0.5cm]
    \textbf{\Large 93044 Regensburg} }
\opening{ Sehr geehrte Damen und Herren,}

hiermit lege ich fristgerecht Widerspruch gegen die Bescheinigung vom
12.05.2026 ein.

Die Bescheinigung wird dem gestellten Antrag in der Sache nicht vollständig
gerecht. Beantragt war die Feststellung eines (steuerlich relevanten)
Grundlagenbescheids zur Behinderung meines Sohnes Leopold, insbesondere die
Feststellung des Grades der Behinderung seit Geburt. Die nun erteilte
Bescheinigung weist demgegenüber einen Zeitraum von 01.09.2024 bis 16.12.2024
aus und ist ausdrücklich als „nicht für steuerliche Zwecke nutzbar“ bezeichnet.
Damit wird das maßgebliche Begehren, nämlich die Feststellung der Behinderung
ab Geburt und in einer für steuerliche Zwecke verwertbaren Form, nicht erfüllt.

Soweit in der Bescheinigung darüber hinaus die gesundheitlichen Voraussetzungen
für die Merkzeichen G, B, aG und H festgestellt werden, war dies nicht die
Intention meines ursprünglichen Antrags. Im Vordergrund stand die Frage der
erstmaligen und zeitlich zutreffenden Feststellung der Behinderung. Die Frage
etwaiger Merkzeichen, insbesondere aG, bedarf einer gesonderten Würdigung und
soll hier nicht im Mittelpunkt stehen. 

Aus meiner Sicht ist entscheidend, dass die gesundheitliche Beeinträchtigung
meines Sohnes bereits kurz nach der Geburt dokumentiert und erkennbar war. Im
Entlassbrief vom 01.06.2022 wird bereits auf einen möglichen Gendefekt
hingewiesen, es wird ausdrücklich das Silver-Russell-Syndrom erwähnt. Außerdem
wird geraten, eine genetische Abklärung durchzuführen. Den Brief finden Sie
anbei. Damit bestanden schon früh nach der Geburt konkrete medizinische
Hinweise auf eine erhebliche gesundheitliche Beeinträchtigung. 

Vor diesem Hintergrund erscheint eine Feststellung der Behinderung erst ab
September 2024 nicht ausreichend. Weitere Arztberichte und Gutachten kann ich
auf Anfrage gerne vorlegen, das Gutachten zum Gendefekt meines Sohnes mit der
Diagnose partielle Trisomie 7 liegt Ihnen bereits vor.  

Ich bitte daher um erneute Prüfung und um eine Abänderung der
Bescheinigung dahingehend, dass der Grad der Behinderung ab Geburt bzw. ab dem
frühestmöglichen Zeitpunkt festgestellt und ein für steuerliche Zwecke
verwendbarer Grundlagenbescheid erteilt wird.  

Hilfsweise beantrage ich die erneute Entscheidung unter vollständiger
Berücksichtigung des Entlassbriefs vom 01.06.2022 sowie der dort bereits
dokumentierten frühzeitigen Hinweise auf einen Gendefekt und das
Silver-Russell-Syndrom.

\closing{Mit freundlichen Grüßen} 
\end{letter}
\end{document}
